\section{単回帰分析と因果推論の基礎}
\label{b07_Regression}

\subsection{授業内容}

\subsubsection{概要}
前回学んだ相関係数は変数間の関連性を示すものであるが，一方の変数から他方の変数を予測・説明する枠組みとして回帰分析がある。本コマでは最も基本的な単回帰分析について学び，予測モデルの構築と評価の方法を理解する。また，相関関係と因果関係の違いについて考察し，因果推論の基礎的な考え方を学ぶ。古典的テスト理論から測定に含まれる誤差を知り，誤差が正規分布に従っていると考えることで，平均の位置を推定することができるようになった考え方を拡張し，他の変数との関係から平均の位置が関数で表現される最も単純な場合である線形モデルについて考える。

\subsubsection{コマ主題細目}
\begin{description}
	\item[予測と線形モデル] 変数間関係を記述する最も単純な形として，一次関数を用いたものを考える。現象に対して関数で記述する数学的現象主義，変数を説明変数と被説明変数として捉えること，予測に誤差が加わって現象となることなど，これまでの議論をデータと数式に置き換えただけであることをしっかりと理解する。さらに単回帰分析の数式的記述について理解する。とくにデータ($X_i,Y_i$)，予測値($\hat{Y}$)，誤差($e_i$)，回帰係数($b_0,b_1$)として表される記号の細部にまで注意を払い，その意味を理解する。
	      \begin{flushright}
		      $\to$ \textcite{Toyoda201706}のP.27--28.
	      \end{flushright}

	\item[最小二乗法] 最小二乗法によって推定値が算出できることを理解する。数式の細かい展開はここでは行わず，最小化する目的関数の直観的理解に止める。また算出された係数(傾きと切片)がどのような意味を持つのか改めて理解する。傾きは説明変数が1単位変化したときの被説明変数の変化量を表し，切片は説明変数がゼロのときの被説明変数の値を表す。これらの係数の解釈と実践的な意味について考察する。
	      \begin{flushright}
		      $\to$　\textcite{Toyoda201706}のP.30
	      \end{flushright}

	\item[残差と決定係数] 回帰係数が計算されれば，実際のデータに対して予測値を計算することができ，予測がどの程度適合していたかを検証することができる。予測値と被説明変数との差分を残差(residuals)と呼ぶが，この両者の相関係数を重相関係数，その二乗したものを決定係数($R^2$)と呼んでモデルの適合度を測る指標にすることを理解する。決定係数は被説明変数の分散のうち，モデルによって説明される割合を表し，モデルの説明力を示す重要な指標である。
	      \begin{flushright}
		      $\to$\textcite{Nojima201903}のP.58-60.
	      \end{flushright}

	\item[相関関係と因果関係] 相関関係があるところに，因果的な説明をしてしまいがちであるが，因果的な説明をするためには相関以上に注意が必要である。因果関係を見るために必要な，時間的先行性，他の原因がないこと，他の結果に至らないことなどの諸条件を確認する。回帰分析は変数間の関係を数理モデルとして表現する強力なツールであるが，それだけでは因果関係を示すには不十分であることを理解する。回帰係数が因果効果として解釈できる条件と，そうでない場合の違いについて理解する。
	      \begin{flushright}
		      $\to$ 相関と因果の違いについては，\textcite{Toyoda1992}に簡潔にまとめられている。
	      \end{flushright}

\end{description}
\subsubsection{キーワード}
\begin{itemize}
	\item 一次関数と線形モデル
	\item 予測値，誤差，回帰係数
	\item 最小二乗法と決定係数
	\item 相関関係と因果関係
\end{itemize}
\subsection{授業情報}
\paragraph{コマの展開方法}
講義と演習
\paragraph{標準シラバスにおける位置づけ}
科目番号5 心理学統計法；1.心理学で用いられる統計手法；C 記述統計：回帰
\subsubsection{予習・復習課題}
\paragraph{予習}
前回学習した相関係数の概念と計算方法について復習しておくこと。また，一次関数の基本的な性質(傾きと切片の意味)についても確認しておくと良い。
\paragraph{復習}
\textcite{Yamauchi201001}のP.62--81は相関係数と回帰分析を1つの章にまとめて解説しているので参考にすると良い。最小二乗法による回帰係数の算出については，偏微分連立方程式を用いれば簡単に計算できるが，偏微分連立方程式とは何かに遡って解説している本として\textcite{kosugi2018}のP.104--112がある。

また，授業で学んだ内容を応用して以下の課題に取り組むこと：
\begin{enumerate}
	\item 実際のデータセットを用いて単回帰分析を実施し，回帰係数と決定係数を算出する
	\item 得られた回帰式を用いて予測を行い，残差をプロットする
	\item 身近な例から「相関関係と因果関係の混同」の事例を見つけ，なぜそれが因果関係とは言えないのかを説明する
	\item 因果関係を主張するためには，相関関係に加えてどのような条件や証拠が必要かを考察する
\end{enumerate}
